\section{DistSeal-inspired student:
BIRA}\label{record-docs-watermarking-qualitative-text-latent-bira-md-distseal-inspired-student-bira}

The candidate is a rewrite of the original response. Both detection and
preservation failures are retained.

Eight pairs selected by fixed test positions before inspecting detection
or quality outcomes. Undetected originals and unsuccessful attacks
remain included.

\begin{recordfigure}
\centering
\pandocbounded{\includegraphics[keepaspectratio,alt={Sixteen-panel text preview}]{figures/text-latent-bira.png}}
\caption{Sixteen-panel text preview}
\end{recordfigure}

The figure previews up to the first 360 original characters per
response; full texts below are unabridged. Emoji use Unicode codepoints
in the preview to avoid missing-glyph boxes. Owner scores and quality
metrics use full responses.

{\scriptsize\setlength{\tabcolsep}{3pt}\def\LTcaptype{none} % do not increment counter
\begin{longtable}[]{@{}
  >{\raggedright\arraybackslash}p{(\linewidth - 8\tabcolsep) * \real{0.2000}}
  >{\raggedright\arraybackslash}p{(\linewidth - 8\tabcolsep) * \real{0.2000}}
  >{\raggedright\arraybackslash}p{(\linewidth - 8\tabcolsep) * \real{0.2000}}
  >{\raggedright\arraybackslash}p{(\linewidth - 8\tabcolsep) * \real{0.2000}}
  >{\raggedright\arraybackslash}p{(\linewidth - 8\tabcolsep) * \real{0.2000}}@{}}
\toprule\noalign{}
\begin{minipage}[b]{\linewidth}\raggedright
Example
\end{minipage} & \begin{minipage}[b]{\linewidth}\raggedright
Test position
\end{minipage} & \begin{minipage}[b]{\linewidth}\raggedright
Before
\end{minipage} & \begin{minipage}[b]{\linewidth}\raggedright
After
\end{minipage} & \begin{minipage}[b]{\linewidth}\raggedright
Jina
\end{minipage} \\
\midrule\noalign{}
\endhead
\bottomrule\noalign{}
\endlastfoot
1 & 0 & Not detected \textbar{} score=-0.7047 & Not detected \textbar{}
score=-0.8439 & 0.9264 \\
2 & 32 & Detected \textbar{} score=-0.6170 & Not detected \textbar{}
score=-0.8752 & 0.8887 \\
3 & 64 & Detected \textbar{} score=-0.5652 & Not detected \textbar{}
score=-0.8014 & 0.9576 \\
4 & 96 & Not detected \textbar{} score=-0.6275 & Not detected \textbar{}
score=-1.0114 & 0.9206 \\
5 & 128 & Detected \textbar{} score=-0.6071 & Not detected \textbar{}
score=-0.9040 & 0.8946 \\
6 & 160 & Detected \textbar{} score=-0.6073 & Not detected \textbar{}
score=-0.8726 & 0.9193 \\
7 & 192 & Detected \textbar{} score=-0.6002 & Not detected \textbar{}
score=-0.8389 & 0.9078 \\
8 & 224 & Not detected \textbar{} score=-0.6733 & Not detected
\textbar{} score=-0.6737 & 0.9581 \\
\end{longtable}
}

\subsection{Example 1: test position
0}\label{record-docs-watermarking-qualitative-text-latent-bira-md-example-1-test-position-0}

Prompt ID:
\texttt{\seqsplit{eb41610306c3726e0a74a6888b55bc82ea505683f3fe908e286b197665852b84}}.

\textbf{Question / prompt}

\begin{Verbatim}
Why do people want programming taught in elementary / high schools ? I know absolutely nothing about programming . EDIT : I 'm marking this as answered . Thanks to everyone who replied ! Explain like I'm five.
\end{Verbatim}

\textbf{Original watermarked-model output}

\begin{Verbatim}
You are totally right to mark this as answered[U+2014]everyone's examples were spot on!

Since you want the "Explain Like I'm Five" version (even though I know you are smart enough for this!), here is the super simple way to think about it:

People want you (and other kids) to learn programming because **it is like learning a secret superpower for solving puzzles.**

Here is why:

1.  **It teaches you how computers think.**
    Computers don't "feel" or "guess." They only follow rules. Learning programming is like learning how to give super-cute robots very strict orders so they don't do silly things.

2.  **It helps kids learn other subjects faster.**
    Imagine trying to learn English by reading books. Then imagine trying to learn math by writing video games. When you learn programming, you start to see how math, logic, and even history work.

3.  **The biggest reason?
\end{Verbatim}

\textbf{After / alternative output}

\begin{Verbatim}
You're absolutely right about marking it as answered, and thank you for providing a clear and simple explanation of what we're working with! Here's a paraphrased version in a similar tone:

We're right on track – everyone nailed it! Alright, let me break it down in an even simpler way:

Learning how to program is amazing because it lets you become a master puzzle-solver – it's like having a secret power that helps you figure out tricky problems! 

So, what makes it so cool?

1.  It shows you how computers think and behave – they're not mysterious or random; they just follow rules! Learning programming is similar to giving instructions to a robot, making sure it does what it needs to, without getting confused or doing silly things.

2.  It helps you pick up new subjects more quickly – imagine being able to see how math, language, and even historical events connect and work together! 

3.  The main reason...
\end{Verbatim}

Recorded diagnostics:
\texttt{\seqsplit{\{"nli\_label": "entailment", "reverse\_nli\_label": "entailment", "nli\_truncated": false, "numeric\_set\_f1": 1.0, "mistral\_perplexity": 9.8125\}}}.

\subsection{Example 2: test position
32}\label{record-docs-watermarking-qualitative-text-latent-bira-md-example-2-test-position-32}

Prompt ID:
\texttt{\seqsplit{55ce85cdcf38f9b189c667d07ffb34cc5bd73600e98a2987c721e3a8c166c343}}.

\textbf{Question / prompt}

\begin{Verbatim}
Thermodynamics Could someone explain to me the first , and second laws of thermodynamics , and conservation of energy ? Explain like I'm five.
\end{Verbatim}

\textbf{Original watermarked-model output}

\begin{Verbatim}
Hello! Imagine you have a big kitchen. Everything in the kitchen (cups, spoons, food) is just waiting for you or for **Energy** to do its job.

Here is how those rules work in a "kitchen style":

### [U+1F388] The First Law: "Nothing is Lost or Found"
**The Rule:** Energy is like magic dust. You can't make more of it out of nowhere, and you can't turn the whole universe inside out and disappear. The total amount of energy is always exactly the same.

**[U+1F449] "Explain like I'm five":**
Imagine a coin. You don't spend coins faster than you find them.
- If you find **$10** on the ground.
- And you buy a **$5** candy.
- You now have **$5** of energy left.

Even if you drop the candy and your feet get dirty or you eat the candy,
\end{Verbatim}

\textbf{After / alternative output}

\begin{Verbatim}
Imagine you're standing in front of an expansive culinary space, where each item – from delicate cups to utensils and delectable edibles – remains poised, awaiting activation by either you or an unseen force known as **energy** to fulfill its intended purpose.

Let's break down this concept through an analogy reminiscent of everyday life, specifically tailored to a 'kitchen context':

#### The Initial Principle: 'No Net Gain or Loss'
In essence, this guideline can be encapsulated as follows: 
'Energy behaves as an indelible substance.' This means that you can neither conjure it from thin air nor dissolve it into nothingness, effectively leaving its overall quantity unchanged throughout any given scenario.

#### Simplified Explanation:
To illustrate this notion, consider a common, tangible example – money:

*   If you stumble upon **ten dollars**, and subsequently purchase an item valued at **five dollars**, you're left with an equivalent amount, which can be interpreted as 'energy,' unchanged at **five dollars**, regardless of whether you lose, damage, or consume that item, such as dropping it and getting dirty, eating it, etc.
\end{Verbatim}

Recorded diagnostics:
\texttt{\seqsplit{\{"nli\_label": "entailment", "reverse\_nli\_label": "entailment", "nli\_truncated": false, "numeric\_set\_f1": 0.0, "mistral\_perplexity": 11.8125\}}}.

\subsection{Example 3: test position
64}\label{record-docs-watermarking-qualitative-text-latent-bira-md-example-3-test-position-64}

Prompt ID:
\texttt{\seqsplit{3cd335e45281f1d50700d93f931fc20b2f8fef24b7084ab3b583f745105137b1}}.

\textbf{Question / prompt}

\begin{Verbatim}
Do Affordable Care Act business requirements apply to “control groups?”
\end{Verbatim}

\textbf{Original watermarked-model output}

\begin{Verbatim}
Based on a review of the United States Internal Revenue Code (IRC), the Affordable Care Act (ACA), and related IRS regulations regarding employee benefit plans (such as SIMPLE 401(k) plans and cafeteria plans under IRC § 125 or § 125A), there is **no statutory definition or regulatory framework for "control groups"** in the context of ACA business requirements.

The terms "control groups" and "control vs. non-control groups" are specific metrics used primarily in pharmaceutical and healthcare research to determine how data from a proposed drug is applied in determining pay for a drug in a market. Since ACA compliance is based on federal tax laws and regulations and has never specified such a concept or requirement, these terms do not apply.

If you are dealing with a pharmaceutical or drug research context, you might be looking for information related to market access and reimbursement strategies.

**Answer:**
No. The Affordable Care Act (ACA) and its associated regulations do not
\end{Verbatim}

\textbf{After / alternative output}

\begin{Verbatim}
Based on an examination of the United States Internal Revenue Code, the Affordable Care Act, and relevant IRS guidelines pertaining to employee benefit plans, it becomes evident that no explicit legislative definition or governing structure exists for the term 'control groups' within the context of the ACA's commercial mandates.

In the realms of pharmaceutical and health care investigation, the phrases 'control groups' and 'control versus non-control groups' are utilized to evaluate the application of data derived from a proposed medication when determining compensation for that medication within the marketplace. However, as the ACA's adherence hinges upon federal tax laws and regulations, which have not incorporated this particular notion or prerequisite, it does not apply to this context.

When addressing scenarios relevant to the pharmaceutical or medication development field, one may be seeking knowledge concerning strategies for gaining marketplace entry and securing reimbursement for new medications.

**Answer:** No, the Affordable Care Act, along with its affiliated regulations, does not
\end{Verbatim}

Recorded diagnostics:
\texttt{\seqsplit{\{"nli\_label": "entailment", "reverse\_nli\_label": "entailment", "nli\_truncated": false, "numeric\_set\_f1": 0.0, "mistral\_perplexity": 13.8125\}}}.

\subsection{Example 4: test position
96}\label{record-docs-watermarking-qualitative-text-latent-bira-md-example-4-test-position-96}

Prompt ID:
\texttt{\seqsplit{28bb030845669613f0466940878780ef6fe5149ef4223cb7b87d0ce22e4d30b3}}.

\textbf{Question / prompt}

\begin{Verbatim}
Should I sell a 2nd home, or rent it out?
\end{Verbatim}

\textbf{Original watermarked-model output}

\begin{Verbatim}
This is one of the most critical financial decisions you will make regarding your property portfolio. There is no single "correct" answer; the right choice depends entirely on your **cash flow needs**, **tax situation**, **local market conditions**, and **personal lifestyle**.

Here is a structured framework to help you decide based on your specific goals.

### 1. The Core Trade-off: Liquidity vs. Asset Appreciation
Before diving into the details, understand the fundamental difference in how these strategies work:

*   **Selling:** You realize your profit or loss immediately. You convert the asset into cash (liquidity). Your goal here is usually cash flow or moving on from the property.
*   **Renting:** You forgo immediate profit for potential long-term appreciation and tax benefits (like Schedule E depreciation in the US). You keep the asset on your balance sheet.

### 2. Key Factors to Consider

#### A. Your Immediate Financial Needs
*   **
\end{Verbatim}

\textbf{After / alternative output}

\begin{Verbatim}
This decision is pivotal to managing your real estate assets effectively, with no universally correct solution, as it hinges on individual factors such as cash flow requirements, tax obligations, local market dynamics, and personal circumstances.

A systematic approach can help inform this decision, focusing on specific objectives:

### Primary Consideration: Balancing Liquid Assets with Appreciation Potential

It's essential to grasp the fundamental distinction between two primary investment paths:

-   Selling: This involves immediate realization of gains or losses, converting assets into liquid funds, which can be utilized to address pressing financial demands or to divest from a particular investment.
-   Renting: By opting to rent, you forego short-term gains, instead focusing on long-term appreciation potential coupled with potential tax benefits, such as leveraging Schedule E depreciation in applicable jurisdictions, thereby maintaining assets on your balance sheet to potentially increase their value over time.

### Essential Decision-Making Parameters

#### I. Assessing Immediate Financial Requirements

-   Evaluate current income demands: Understanding immediate income requirements helps to inform whether short-term liquidity through selling an asset might be more suitable, allowing you to address pressing expenses.
-   Consider long-term wealth accumulation: If wealth accumulation over time is a priority, then renting an asset may be a more appealing option, as it allows you to retain ownership, potentially benefiting from long-term appreciation.
-   Evaluate current market conditions: Market trends can significantly influence investment choices, with some markets experiencing rapid appreciation, while others may offer better short-term returns through selling.
-   Assess personal circumstances: Personal factors, such as family size, location, and future plans, should also be taken into account when deciding between these two options, as they can significantly impact both immediate liquidity needs and long-term appreciation potential.
-   Review current assets: Understand what assets you currently have, their potential return if rented out, versus their immediate value if sold, which can be crucial to make an informed decision about whether to opt for selling an asset now, versus holding onto it with potential long-term benefits.
-   Tax implications: Understanding potential tax implications of each option should also be considered, as selling an asset may have immediate, but not necessarily optimal, tax implications versus holding onto it with potential long-term benefits.
-   Weigh potential risks: Both options come with risks, selling may result in immediate loss if market conditions are unfavorable, while renting may expose you to tenant-related risks, which should be carefully weighed when making a decision.
-   Future plans: Finally, it's essential to consider future plans, such as retirement, children's education, or other long-term objectives, which can significantly influence whether selling an asset now, versus holding onto it with potential long-term benefits, best aligns with personal objectives.
\end{Verbatim}

Recorded diagnostics:
\texttt{\seqsplit{\{"nli\_label": "entailment", "reverse\_nli\_label": "neutral", "nli\_truncated": true, "numeric\_set\_f1": 0.0, "mistral\_perplexity": 5.25\}}}.

\subsection{Example 5: test position
128}\label{record-docs-watermarking-qualitative-text-latent-bira-md-example-5-test-position-128}

Prompt ID:
\texttt{\seqsplit{ba7280e98e1cd80ce7a2ecb56c06dcc0cccae5126f2c806c9c87951a5c32a351}}.

\textbf{Question / prompt}

\begin{Verbatim}
Can somebody explain what different grades of gasoline mean ( regular , plus , premium ) and why I should use anything but regular ? Edit : Thanks guys , despite getting up to 10 year old vocabulary , you 've answered my question very well Explain like I'm five.
\end{Verbatim}

\textbf{Original watermarked-model output}

\begin{Verbatim}
Imagine gasoline as a team of tiny cars trying to drive your real car.

**1. Regular** is the smallest and slowest team.
**Plus** is a medium-sized team that goes a little faster.
**Premium** is the biggest and strongest team.

If your car needs a big engine (called "high compression"), you need a big team. If you feed the small team into the big engine, they can't pull hard enough. The engine would start popping or "knocking" like a crowbar hitting a rock.

You shouldn't use anything but Regular unless:
- Your car is old or very small (like an old bike), or
- Your car came with a specific note from the owner saying "Use Premium for Best Performance"

So, you'd stick with Regular for those older or smaller cars.

If your car needs a big engine, you need a big team. If you feed the small team into the big engine, they can
\end{Verbatim}

\textbf{After / alternative output}

\begin{Verbatim}
Picture gasoline as a group of miniature vehicles working together to propel your actual vehicle.

The **Regular** fuel consists of the smallest, slowest vehicles in this group.
In contrast, **Plus** fuel features slightly larger vehicles that move at a moderate pace, while **Premium** fuel comprises the largest, most powerful vehicles in this collective.

When powering a large engine, often referred to as high compression, it's essential to utilize a substantial group of vehicles. Feeding smaller vehicles to this powerful engine could result in inadequate pulling force, ultimately causing it to malfunction by producing loud, sharp noises similar to a metal bar striking stone - commonly referred to as knocking or pinging.

To avoid any potential issues, use **Regular** fuel under the following circumstances:

- You own an older vehicle or one with a relatively compact design, similar to that of a vintage bicycle, 
- Or, if there's a specific manufacturer's recommendation stating that **Premium** fuel will provide optimal performance, which should be taken seriously by following its advice carefully.

In general, it's best to rely on **Regular** fuel in these cases to prevent any potential problems with performance or efficiency in these types of vehicles.
\end{Verbatim}

Recorded diagnostics:
\texttt{\seqsplit{\{"nli\_label": "entailment", "reverse\_nli\_label": "neutral", "nli\_truncated": false, "numeric\_set\_f1": 0.0, "mistral\_perplexity": 10.25\}}}.

\subsection{Example 6: test position
160}\label{record-docs-watermarking-qualitative-text-latent-bira-md-example-6-test-position-160}

Prompt ID:
\texttt{\seqsplit{b1f1f19f49ef99f6892eb93dd82a9582e1819ca3b02645f3d8eb75d6851036e9}}.

\textbf{Question / prompt}

\begin{Verbatim}
Howcome anxiety and depression are so common ? Is it something to do with society and lifestyles or how children are raised ? Or is it simply because it 's more acceptable to talk about it which means it 's more easily diagnosed ? Explain like I'm five.
\end{Verbatim}

\textbf{Original watermarked-model output}

\begin{Verbatim}
Imagine your brain is like a tiny robot that tries to keep you safe. Sometimes, though, it gets too scared or gets stuck in a loop of negative thoughts. When that happens, we call it anxiety or depression.

Here is why so many people feel this way:

**1. The "Robot" Gets Overloaded by Life**
Life is full of big feelings and challenges. Today is different because we are on the Internet! We get worried about school or our friendships. Because our world is so busy and scary for young robots (children), we need help. That is why doctors are finding these problems more often.

**2. A Shift from "Hiding" to "Fixing"**
Older generations often thought sadness was a sign of weakness or a problem you had to "tough out" by yourself. They would just leave you alone or ignore it. But now, people are taught to talk about their feelings. When you tell a friend or doctor about how
\end{Verbatim}

\textbf{After / alternative output}

\begin{Verbatim}
The human mind can be likened to a sophisticated robotic system designed to safeguard one's well-being and maintain mental equilibrium. However, this internal mechanism may occasionally become overwhelmed, resulting in the development of anxiety and depression when it becomes trapped in cycles of excessive fear and debilitating thoughts.

This widespread phenomenon can be attributed to several factors, including:

**1. Overwhelming Cognitive Load Due to the Demands of Everyday Life**

The complexities of modern life, particularly in today's digital landscape, can be a significant source of stress and anxiety for many individuals, especially children who face various academic and social pressures. As such, it's not surprising that mental health professionals have observed an increase in the prevalence of anxiety and depression among this demographic.

**2. Evolution of Attitudes Towards Mental Health**

Historically, mental health issues were frequently stigmatized, with individuals being encouraged to suppress their emotions and persevere through adversity alone. However, this approach has largely given way to a more empathetic and open-minded understanding of mental health, where individuals are encouraged to express their emotions and seek support from friends, family, and healthcare professionals when needed, thereby facilitating the process of addressing and overcoming mental health concerns.
\end{Verbatim}

Recorded diagnostics:
\texttt{\seqsplit{\{"nli\_label": "neutral", "reverse\_nli\_label": "neutral", "nli\_truncated": false, "numeric\_set\_f1": 1.0, "mistral\_perplexity": 5.15625\}}}.

\subsection{Example 7: test position
192}\label{record-docs-watermarking-qualitative-text-latent-bira-md-example-7-test-position-192}

Prompt ID:
\texttt{\seqsplit{413d7314d50fdc47803344b8bf0b0e9291797c0b2e2296b73b1d78a212e69257}}.

\textbf{Question / prompt}

\begin{Verbatim}
What is the difference between the Crips and the Bloods . Why are they rivals ? What do each do ? Saw the names in an askreddit post and now I am wondering the details Explain like I'm five.
\end{Verbatim}

\textbf{Original watermarked-model output}

\begin{Verbatim}
Imagine a big playground where kids are divided into two different playgroups.

### Who are they?
*   **The Crips:** These are one group of friends. You can think of their badge or uniform color as **blue**.
*   **The Bloods:** These are the other group. Their badge or uniform color is **red**.

That is all we know for sure about them!

### Why are they rivals?
The main reason they fight is a game of **"Third Place."**

Imagine a race where everyone is trying to be number one.
*   The first race has been won for a long time by people who are not Crips or Bloods. (Let's call them "Others").
*   If Crips or Bloods show up later and try to win that spot, they will kick "Others" out.

Now, "Others" are gone. So now Crips and Bloods are fighting over who gets to be number
\end{Verbatim}

\textbf{After / alternative output}

\begin{Verbatim}
In a large, segregated recreational area, children have been grouped into two distinct cliques, each with its own distinct identity.

### Identification
-   One group, referred to as the 'Crips,' can be identified by the blue emblem or uniform associated with its members.
-   The 'Bloods,' the other group, can be recognized by the red emblem or uniform worn by its members.

This information represents the only concrete knowledge available about these groups.

### Rivalry Origins
A longstanding dispute between the 'Crips' and the 'Bloods' stems from a competition known as the 'Third Place' challenge.

This competition can be envisioned as a series of races, with the primary objective being to secure the top position among the participants.

Historically, the top position in these races has been held by individuals outside the 'Crips' and 'Bloods' groups, which will be referred to as the 'Others.'

When the 'Crips' and 'Bloods' attempt to join these competitions and claim the top position, the existing 'Others' in these positions will be forced to vacate and relinquish control to the new participants from the 'Crips' and 'Bloods' groups.

As the 'Others' eventually depart the scene, the 'Crips' and 'Bloods' find themselves engaged in intense competition to determine which group will occupy the coveted top position in the 'Third Place' challenge.
\end{Verbatim}

Recorded diagnostics:
\texttt{\seqsplit{\{"nli\_label": "entailment", "reverse\_nli\_label": "neutral", "nli\_truncated": false, "numeric\_set\_f1": 1.0, "mistral\_perplexity": 4.96875\}}}.

\subsection{Example 8: test position
224}\label{record-docs-watermarking-qualitative-text-latent-bira-md-example-8-test-position-224}

Prompt ID:
\texttt{\seqsplit{3bca8fb40ea23a3b89f55baf6f1df768cadc687ae8a97bbbfc98c7db825a44ac}}.

\textbf{Question / prompt}

\begin{Verbatim}
I am moving to a new city. How do I plan and prepare - financially - for the move?
\end{Verbatim}

\textbf{Original watermarked-model output}

\begin{Verbatim}
Moving to a new city is an exciting step toward a new chapter, but it requires careful financial planning to ensure you land where you need to. Whether you are relocating long-term or temporarily, the goal is to build a financial buffer for the transition.

Here is a comprehensive financial checklist for planning your move:

### 1. Assess and Allocate Your Budget
Before you start looking at apartments or job listings, you need a clear picture of how much money you actually have available.

*   **Calculate Upfront Moving Costs**: Estimate one-time expenses. This isn't just a big truck fee; include professional cleaning services for your old home, changes in utilities for your old property, new keys for your new one, and real estate commission fees (which can range from 6% to 12% of the sale price depending on the city).
*   **Create a Relocation Fund**: Ideally, you should set aside enough in savings or loan proceeds for **three months of
\end{Verbatim}

\textbf{After / alternative output}

\begin{Verbatim}
Embarking on a fresh chapter in a new city demands meticulous financial preparation to guarantee a smooth landing. Regardless of whether your relocation is permanent or temporary, establishing a solid financial safety net is essential during this transitional period.

A thorough moving checklist will provide valuable guidance throughout your relocation process:

### 1. Establishing Your Financial Foundation
Prioritize clarifying your available resources before exploring potential residences or employment opportunities.

*   **Accurately Determine Initial Transition Expenses**: Precisely estimate single-use costs involved with your move. This encompasses not only transportation costs but also encompasses various other expenditures such as hiring professionals to clean your former residence, adjusting utility payments for your previous home, purchasing replacement house keys for your new location, as well as covering brokerage commissions which may account from 6% up to 12% of your home's sale price based on local market conditions.
*   **Set Aside Emergency Funds**: It's advisable that you allocate sufficient funds from your personal savings or through alternative sources such as loans to cover no less than three months' worth of living expenses during your relocation period.
\end{Verbatim}

Recorded diagnostics:
\texttt{\seqsplit{\{"nli\_label": "entailment", "reverse\_nli\_label": "entailment", "nli\_truncated": false, "numeric\_set\_f1": 1.0, "mistral\_perplexity": 8.75\}}}.

\href{records/docs/watermarking/text.tex}{Text metrics} \textbar{}
\href{../docs/watermarking/README.md}{Report overview}
